\documentclass[12pt]{article}
\usepackage[utf8]{inputenc}
\usepackage[T1]{fontenc}
\usepackage{hyperref}
\title{Polabian Language Resources for K-12 Curriculum Development}
\author{Compiled by Claude Sonnet 5.5 (claude-sonnet-5-5) for Andr\'as Kornai}
\date{}
\begin{document}
\maketitle

\section{Introduction}
Polabian (Drevanian-Polabian, Lüneburg Wendish; ISO 639-3 pox) is an extinct West Slavic language of northern Germany; its last known speaker is reported to have died in 1756. There are no native speakers, and the language is attested only through short word lists, glossaries and a few texts recorded by 17th and 18th century writers. Any curriculum would therefore be a heritage and reconstruction programme built on scholarly sources.

\section{Primary Educational Resources}
In the searches made for this report I found no school textbooks, primers or classroom programmes for Polabian.

\section{Linguistic and Reference Materials}
\begin{itemize}
\item Overview of the language and its sources: \url{https://en.wikipedia.org/wiki/Polabian_language}
\item Wiktionary entries for Polabian (Drevani) words, which cite the Polish etymological dictionary and the Polański--Sehnert dictionary: \url{https://en.wiktionary.org/wiki/c%C3%BC}
\item A user page on German Wikipedia listing sources for a planned Polabian dictionary: \url{https://de.wikipedia.org/wiki/Benutzer:Phillipm0703/Polabisches_W%C3%B6rterbuch/Quellen}
\item Scholarly PDFs from the Düsseldorf university repository that appeared in search results (content not verified in detail): \url{https://lessing.ub.uni-duesseldorf.de/download/pdf/5989133.pdf} and \url{https://lessing.ub.uni-duesseldorf.de/download/pdf/7296464.pdf}
\end{itemize}
The search results also mention, without a freely accessible URL, the Polański and Sehnert \emph{Polabian-English Dictionary} (1967), the Lehr-Spławiński and Polański etymological dictionary, and Olesch's \emph{Thesaurus Linguae Dravaenopolabicae}.

\section{Religious and Community Resources}
No religious or community resources were found.

\section{Digital and Audio Resources}
No audio or dedicated digital learning resources were found; Wiktionary is the only online lexical resource seen.

\section{Curriculum Development Context}
Because the language is extinct and sparsely attested, material would have to be adapted from the scholarly dictionaries and word lists above, and any use of the language in school would be limited to reconstructed vocabulary and short texts.
\end{document}